% Generated draft: requires mathematical and source review.
\documentclass[11pt]{article}
\usepackage[margin=25mm]{geometry}
\usepackage{fontspec}
\setmainfont{Latin Modern Roman}
\usepackage{amsmath,amssymb,mathtools,unicode-math}
\setmathfont{Latin Modern Math}
\usepackage{enumitem,needspace,xurl,hyperref}
\setlength{\parindent}{0pt}
\setlength{\parskip}{5pt}
\setlength{\emergencystretch}{3em}
\newcommand{\N}{\mathbb{N}}
\newcommand{\Z}{\mathbb{Z}}
\newcommand{\Q}{\mathbb{Q}}
\newcommand{\R}{\mathbb{R}}
\newcommand{\C}{\mathbb{C}}
\newcommand{\F}{\mathbb{F}}
\newcommand{\A}{\mathbb{A}}
\newcommand{\PP}{\mathbb{P}}
\newcommand{\E}{\mathbb{E}}
\newcommand{\eps}{\varepsilon}
\newcommand{\dd}{\,\mathrm{d}}
\begin{document}
\section*{UnsolvedMath ID: 30003644}
\textbf{Source problem code:} catalogue entry
\subsection{Definitions and mathematical statement}
For \(s\in\C\), define the Dirichlet polynomial
\[
D(s)=1+2^{-s}+3^{-s}+5^{-s},\qquad n^{-s}=\exp(-s\log n),
\]
where \(\log n\) is the real natural logarithm for \(n\in\{2,3,5\}\).
The nonvanishing question is whether
\[
\forall s\in\C,\qquad \operatorname{Re}(s)=1\ \Longrightarrow\ D(s)\ne0.
\]
Equivalently, writing \(s=1+it\), it asks whether
\[
\forall t\in\R,\qquad
1+\frac{e^{-it\log 2}}{2}+\frac{e^{-it\log 3}}{3}+\frac{e^{-it\log 5}}{5}\ne0.
\]
The original source record also presents the following as an equivalent Tauberian formulation. Is it true that every function \(f:(0,\infty)\to(0,\infty)\) that is nondecreasing, meaning
\[
0<x\le y\ \Longrightarrow\ f(x)\le f(y),
\]
and satisfies
\[
\lim_{x\to\infty}\frac{f(x)+f(x/2)+f(x/3)+f(x/5)}{x}=\frac{61}{30}
\]
necessarily satisfies
\[
\lim_{x\to\infty}\frac{f(x)}{x}=1?
\]
These limits express the supplied formulas \(f(x)+f(x/2)+f(x/3)+f(x/5)=(61/30+o(1))x\) and \(f(x)=(1+o(1))x\), with the asymptotic direction interpreted as \(x\to\infty\). The original record does not explicitly state that direction. The equivalence of the two questions is reported by the supplied record and is not established in this statement-only draft. No continuity or differentiability assumption on \(f\) is supplied.
\subsection{Short English statement}
The problem asks whether \(1+2^{-s}+3^{-s}+5^{-s}\) is nonzero at every complex point on the line \(\operatorname{Re}(s)=1\). The original record also asks, as a stated equivalent, whether a positive nondecreasing function whose four indicated dilates have sum asymptotic to \((61/30)x\) must itself be asymptotic to \(x\). The supplied open status and verification labels are curation claims requiring independent review.
\subsection{Sources}
\begin{itemize}
\item Local UnsolvedMath registry, supplied complete record for permanent ID \(30003644\), source problem code \texttt{catalogue entry}, title ``Zero-Free Dirichlet Polynomials on the Unit Line.'' Registry collection: \texttt{data/unsolvedmath/problems.json}; supplementary existing-record catalog: \texttt{catalog/problems.json}. This draft uses the supplied record without independently inspecting those files.
\item Supplied primary-source citation: \emph{Mini-Workshop: Interplay between Number Theory and Analysis for Dirichlet Series} (2018), Oberwolfach Reports, pp.~3035--3069, DOI \texttt{10.4171/owr/2017/51}. \url{https://doi.org/10.4171/owr/2017/51}. The record attributes the problem to the workshop's open-problem discussions, gives proposed year 2017, and names no individual proposer.
\item Supplied literature reference: Erd\H{o}s--Ingham, \emph{Arithmetical Tauberian theorems}. \url{https://eudml.org/doc/207489}.
\item Supplied literature reference: Erd\H{o}s Problems, problem 967. \url{https://www.erdosproblems.com/latex/967}.
\item Supplied literature reference: Oberwolfach Report 51, open-problems session. \url{https://doi.org/10.4171/owr/2017/51}.
\item Literature assessment dated exactly \texttt{2026-08-21}: the supplied curation labels the problem open and describes this polynomial as the simplest unresolved case of the Erd\H{o}s--Ingham nonvanishing question on \(\operatorname{Re}(s)=1\). This is supplied dated curation, not independently verified. The background explicitly requests independent verification of time-sensitive status and cited claims.
\item The registry supplies statement status \texttt{corrected\_verified} and states that the clean statement was checked against the official Oberwolfach report and preserves its hypotheses, notation, and open mathematical target. This is a supplied source-verification claim; no source was independently visited for this draft.
\end{itemize}
\end{document}
